\documentclass[a4paper,11pt]{article}

% ---------- Кодировка и язык ----------
\usepackage[utf8]{inputenc}
\usepackage[T2A]{fontenc}
\usepackage[english,russian]{babel}

% ---------- Геометрия страницы ----------
\usepackage[margin=2.5cm]{geometry}

% ---------- Математика ----------
\usepackage{amsmath,amssymb,amsthm,mathtools}
\usepackage{empheq}
\usepackage{bm}          % жирные символы в формулах: \bm{x}

% ---------- Оформление ----------
\usepackage{enumitem}
\usepackage{booktabs}       % красивые таблицы
\usepackage{xcolor}
\usepackage[colorlinks=true,linkcolor=blue,urlcolor=blue,citecolor=blue]{hyperref}
\usepackage{cleveref}       % \cref{...} для умных ссылок
\newcommand{\lean}[1]{\mbox{\path{#1}}} % имена из Lean-кода, без переносов внутри имени (пакет url)

% ---------- Теоремы и окружения ----------
\theoremstyle{plain}
\newtheorem{theorem}{Теорема}[section]
\newtheorem{lemma}[theorem]{Лемма}
\newtheorem{proposition}[theorem]{Утверждение}
\newtheorem{corollary}[theorem]{Следствие}

\theoremstyle{definition}
\newtheorem{definition}[theorem]{Определение}
\newtheorem{example}[theorem]{Пример}

\theoremstyle{remark}
\newtheorem*{remark}{Замечание}

% ---------- Полезные макросы ----------
\newcommand{\R}{\mathbb{R}}
\newcommand{\Z}{\mathbb{Z}}
\newcommand{\N}{\mathbb{N}}
\newcommand{\Q}{\mathbb{Q}}
\newcommand{\norm}[1]{\left\lVert #1 \right\rVert}
\newcommand{\abs}[1]{\left\lvert #1 \right\rvert}
\newcommand{\set}[1]{\left\{ #1 \right\}}
\newcommand{\ip}[2]{\langle #1, #2\rangle} % скалярное произведение
\newcommand{\Rpp}[1][2]{\R^{#1}_{++}} % \Rpp = R^2_{++}, \Rpp[n] = R^n_{++}
\newcommand{\Rp}[1][2]{\R^{#1}_{+}}   % \Rp = R^2_{+}, \Rp[n] = R^n_{+}

\DeclareMathOperator{\conv}{conv}
\DeclareMathOperator{\cone}{cone}
\DeclareMathOperator{\Sp}{Sp}
\DeclareMathOperator{\interior}{int}
\newcommand{\dd}{\widehat d}

\newcommand{\todo}[1]{\textcolor{red}{\textbf{TODO:} #1}}

\title{Наименьшее масштабное возмущение цен\\ для слабой разрешимости}
\date{}

\begin{document}
\maketitle

\section{Введение}

Всюду $d\in\N$, $d\ge2$. Обозначения те же, что в основных заметках. Для $p,x\in\R^d$ через $p\circ x$ обозначено покоординатное произведение, через $\mathbf e:=(1,\dots,1)\in\R^d$~--- единичный вектор, через $\mathbf 1_S$~--- индикатор множества $S$. Неравенства между векторами покоординатные.

Цены~--- наборы из $d+1$ положительных чисел, определённые с точностью до общего положительного множителя. Нормируем их по последней координате и получаем $p\in\Rpp[d]$. Набор цен по периодам есть $P:=\set{p_t\in\Rpp[d]}_{t\in[T]}$, а $y\in\Rp[T]$~--- вектор выпусков.

Из основных заметок используются без доказательства следующие объекты. Вектор $y$ \textit{слабо разрешим} при ценах $P\in\Rp[d\times T]$, если в любой окрестности $P$ найдётся набор $Q\in\Rpp[d\times T]$, при котором $y$ неоклассически разрешим. Множество $S$ называется \textit{слабо достижимым спектром} при ценах $P$, если для каждого $t\in S$ существует ненулевой $\xi\in\Rp[d]$ с $\ip{\xi}{p_s-p_t}\ge0$ при всех $s\notin S$. Семейство таких множеств обозначается $\widehat{\Sp}(P)$. Индексы здесь могут пробегать любое конечное множество точек из $\Rp[d]$.

\begin{theorem}[критерий слабой разрешимости, основные заметки]\label{thm:closure}
    Вектор $y\in\Rp[T]$ слабо разрешим при ценах $P\in\Rp[d\times T]$ тогда и только тогда, когда $y\in\cone\set{\mathbf 1_S\mid S\in\widehat{\Sp}(P)}$.
\end{theorem}

\paragraph{Расстояние.}
Расстоянием между наборами цен $(x_0,\dots,x_d)$ и $(x'_0,\dots,x'_d)$ назовём $\sum_{i=0}^d\abs{\ln x_i-\ln x'_i}$. Для нормированных цен $a,b\in\Rpp[d]$ представляющие наборы имеют вид $(\alpha a,\alpha)$ и $(\beta b,\beta)$. Положим $u_i:=\ln a^i-\ln b^i$ и $s:=\ln\alpha-\ln\beta$. Расстояние между наборами равно $\sum_{i\in[d]}\abs{u_i+s}+\abs{s}$. Минимизируя по $s\in\R$, получаем индуцированное расстояние
\[
    \dd(a,b):=\min_{s\in\R}\Bigl(\sum_{i\in[d]}\abs{u_i+s}+\abs{s}\Bigr).
\]
Минимум достигается, так как функция выпукла, кусочно линейна и стремится к бесконечности. При $d=2$ минимум равен сумме абсолютных отклонений чисел $u_1,u_2,0$ от медианы, то есть разности наибольшего и наименьшего из них:
\[
    \dd(a,b)=\max(\abs{u_1},\abs{u_2},\abs{u_1-u_2}).
\]
Шар $\set{\dd(\cdot,p)\le\rho}$ в логарифмических координатах есть шестиугольник.

\paragraph{Задача.} Если $y$ не слабо разрешим при $P$, ищем ближайший набор $Q$, при котором он слабо разрешим:
\[
    \rho^*:=\inf\Bigl\{\max_{t\in[T]}\dd(q_t,p_t)\;\Big|\;Q\in\Rpp[d\times T],\ y\text{ слабо разрешим при }Q\Bigr\}.
\]

\paragraph{Результаты.}
\begin{enumerate}[label=(\roman*)]
    \item Задача сводится к масштабированию: достаточно искать $q_t=c_tp_t$ с $c_t\in[e^{-\rho},e^{\rho}]$ (теорема~\ref{thm:corners}). Условие записывается как конусное условие на $2T$ угловых точек $e^{-\rho}p_t$, $e^{\rho}p_t$.
    \item Наименьший радиус достигается и является точкой излома (теорема~\ref{thm:least}). Точки излома рациональны как функции от $z=e^{2\rho}$, и их $O\bigl((2T+d)^{d+1}\bigr)$ штук.
    \item Поэтому $\rho^*$ находится бинарным поиском за $O(d\log(T+d))$ обращений к оракулу при фиксированном $d$. Оракул проверяет слабую разрешимость с двусторонними границами. При $d=2$ оракул есть одна задача линейного программирования о потоке (\texttt{poly-hardness/}), и алгоритм полиномиален. При $d\ge3$ оракул NP-труден (\texttt{np-hardness/}), на практике используется MILP.
\end{enumerate}

\section{Метрика}

Пусть $a,b\in\Rpp[d]$ и $u_i:=\ln a^i-\ln b^i$.

\begin{lemma}\label{lem:metric}
    Для всех $i\in[d]$ верно $\dd(a,b)\ge\abs{u_i}$. Для $c>0$ верно $\dd(ca,a)=\abs{\ln c}$.
\end{lemma}

\begin{proof}
    При любом $s$ имеем $\abs{u_i}\le\abs{u_i+s}+\abs{s}\le\sum_j\abs{u_j+s}+\abs{s}$. Для $ca$ и $a$ все $u_i=\ln c$. Значение $s=-\ln c$ даёт $\dd(ca,a)\le\abs{\ln c}$, а первое утверждение даёт $\ge$.
\end{proof}

\begin{corollary}\label{cor:ball}
    Шар $\set{q\in\Rpp[d]\mid\dd(q,p)\le\rho}$ лежит в коробке $e^{-\rho}p\le q\le e^{\rho}p$ и содержит отрезок $\set{cp\mid e^{-\rho}\le c\le e^{\rho}}$.
\end{corollary}

\begin{lemma}[форма шара]\label{lem:shape}
    Пусть $\rho\ge0$. Тогда $\dd(a,b)\le\rho$ тогда и только тогда, когда
    \[
        u\in\rho\conv\set{\pm e_1,\dots,\pm e_d,\pm\mathbf e}.
    \]
\end{lemma}

\begin{proof}
    Функция $\dd$ есть фактор-норма $\ell^1$-нормы на $\R^{d+1}$ по прямой $\R\mathbf e$, вычисленная на векторе $(u,0)$. Её единичный шар есть образ шара $\conv\set{\pm e_1,\dots,\pm e_{d+1}}$ при отображении $(v,v_{d+1})\mapsto v-v_{d+1}\mathbf e$. Минимум по $s$ достигается, поэтому образ замкнутого шара есть замкнутый шар. Отображение переводит $\pm e_i$ в $\pm e_i$ при $i\le d$ и $\pm e_{d+1}$ в $\mp\mathbf e$.
\end{proof}

При $d=2$ это шестиугольник с вершинами $\pm(1,0)$, $\pm(0,1)$, $\pm(1,1)$. При $d=3$ это параллелепипед с вершинами $\pm e_1$, $\pm e_2$, $\pm e_3$, $\pm(1,1,1)$. Вершины $\pm\mathbf e$ лежат на диагонали куба $[-1,1]^d$, и все вершины лежат в этом кубе. Отсюда снова следует следствие~\ref{cor:ball}.

\section{Сведение к масштабам}

Пусть $W=\set{w_r}_{r\in R}\subset\Rpp[d]$~--- конечный набор точек. Для $S\subset R$ положим
\[
    K_S:=\conv\set{w_r\mid r\notin S}+\Rp[d],\qquad K_R:=\varnothing.
\]
Множество $K_S$ выпукло, замкнуто, и $K_S+\Rp[d]\subset K_S$. Для $S\ne R$ оно имеет непустую внутренность и лежит в $\Rpp[d]$.

\begin{lemma}[достижимость и внутренность]\label{lem:interior}
    Пусть $S\subset R$, $t\in S$. Ненулевой $\xi\in\Rp[d]$ с $\ip{\xi}{w_s-w_t}\ge0$ при всех $s\notin S$ существует тогда и только тогда, когда $w_t\notin\interior K_S$. Поэтому $S\in\widehat{\Sp}(W)$ тогда и только тогда, когда ни одна точка $w_t$, $t\in S$, не лежит в $\interior K_S$.
\end{lemma}

\begin{proof}
    Для $\xi\ge0$ условие $\ip{\xi}{w_s-w_t}\ge0$ при $s\notin S$ равносильно $\ip{\xi}{z-w_t}\ge0$ на $K_S$. Если $w_t\in\interior K_S$ и $\xi\ne0$, то точка $w_t-\varepsilon\xi$ лежит в $K_S$ при малом $\varepsilon>0$, и значение равно $-\varepsilon\norm{\xi}^2<0$. Значит, $\xi$ не подходит. Обратно, пусть $w_t\notin\interior K_S$. Если $K_S=\varnothing$, подходит любой ненулевой $\xi\ge0$. Иначе по теореме отделимости точки от открытого выпуклого множества $\interior K_S$ найдётся ненулевая линейная функция $\ip{\xi}{\cdot}$ с $\ip{\xi}{z-w_t}\ge0$ на $\interior K_S$, а значит и на $K_S=\overline{\interior K_S}$. Из $K_S+\Rp[d]\subset K_S$ получаем $\xi\ge0$.
\end{proof}

\begin{lemma}[выход вдоль луча]\label{lem:ray}
    Пусть $S\ne R$ и $p\in\Rpp[d]$. Положим $\gamma:=\min\set{c\ge0\mid cp\in K_S}$. Тогда $\gamma>0$,
    \[
        \set{c\ge0\mid cp\notin\interior K_S}=[0,\gamma],\qquad\set{c\ge0\mid cp\notin K_S}=[0,\gamma).
    \]
    Для $S=R$ положим $\gamma:=+\infty$, и оба множества равны $[0,\infty)$.
\end{lemma}

\begin{proof}
    Минимум существует, так как $K_S$ замкнуто и непусто. Из $K_S\subset\Rpp[d]$ следует $\gamma>0$. Если $cp\in K_S$ и $c'>c$, то $c'p=cp+(c'-c)p$ и $(c'-c)p\in\interior\Rp[d]$, поэтому $c'p\in\interior K_S$. Отсюда $\set{c\mid cp\in K_S}=[\gamma,\infty)$ и $(\gamma,\infty)\subset\set{c\mid cp\in\interior K_S}$. Точка $\gamma p$ не лежит во внутренности, так как $(\gamma-\varepsilon)p\notin K_S$ при $\varepsilon>0$.
\end{proof}

Пусть $W=(a,b)$, $a_t:=e^{-\rho}p_t$, $b_t:=e^{\rho}p_t$, индексы $R=\set{a_t,b_t}_{t\in[T]}$. Для $S\subset R$ и $w:=\sum_km_k\mathbf 1_{S_k}$ значение $w(r)$ есть $\sum_km_k\mathbf 1_{S_k}(r)$.

\begin{theorem}[углы]\label{thm:corners}
    Пусть $d\ge2$, $y\in\Rp[T]$, $\rho\ge0$. Равносильны:
    \begin{enumerate}[label=(\Alph*)]
        \item существует $Q\in\Rpp[d\times T]$ с $\dd(q_t,p_t)\le\rho$, при котором $y$ слабо разрешим,
        \item существуют $c_t\in[e^{-\rho},e^{\rho}]$, при которых $y$ слабо разрешим при ценах $(c_tp_t)_{t\in[T]}$,
        \item существует $w\in\cone\set{\mathbf 1_S\mid S\in\widehat{\Sp}(W)}$ с $w(a_t)\ge y_t\ge w(b_t)$ при всех $t\in[T]$.
    \end{enumerate}
\end{theorem}

\begin{proof}
    (B)$\Rightarrow$(A). По лемме~\ref{lem:metric} $\dd(c_tp_t,p_t)=\abs{\ln c_t}\le\rho$.

    (A)$\Rightarrow$(C). По теореме~\ref{thm:closure} $y=\sum_km_k\mathbf 1_{S_k}$ с $m_k>0$ и $S_k\in\widehat{\Sp}(Q)$. Положим $K_k:=K_{S_k}(Q)$. По следствию~\ref{cor:ball} $a_t\le q_t\le b_t$. Определим $S'_k\subset R$:
    \[
        a_t\in S'_k\iff a_t\notin\interior K_k,\qquad b_t\in S'_k\iff b_t\notin K_k.
    \]
    \emph{Шаг 1: $S'_k\in\widehat{\Sp}(W)$.} Если $a_t\notin S'_k$, то $a_t\in\interior K_k$. Если $b_t\notin S'_k$, то $b_t\in K_k$. Поэтому $K_{S'_k}(W)\subset K_k$, и $\interior K_{S'_k}(W)\subset\interior K_k$. Для $a_t\in S'_k$ имеем $a_t\notin\interior K_k\supset\interior K_{S'_k}(W)$. Для $b_t\in S'_k$ имеем $b_t\notin K_k\supset K_{S'_k}(W)$. По лемме~\ref{lem:interior} $S'_k\in\widehat{\Sp}(W)$.

    \emph{Шаг 2: границы.} Пусть $t\in S_k$. Тогда $q_t\notin\interior K_k$. Если бы $a_t\in\interior K_k$, то $q_t\in a_t+\Rp[d]\subset\interior K_k$. Значит, $a_t\in S'_k$. Пусть $b_t\in S'_k$. Тогда $b_t\notin K_k$, и $q_t\notin K_k$, так как $q_t\le b_t$ и $K_k+\Rp[d]\subset K_k$. Значит, $t\in S_k$. Для $w:=\sum_km_k\mathbf 1_{S'_k}$ получаем $w(a_t)\ge y_t\ge w(b_t)$.

    (C)$\Rightarrow$(B). Пусть $w=\sum_km_k\mathbf 1_{S'_k}$, $m_k>0$, $S'_k\in\widehat{\Sp}(W)$, и $K_k:=K_{S'_k}(W)$. Положим $\gamma_{k,t}:=\min\set{c\ge0\mid cp_t\in K_k}\in(0,\infty]$ и
    \[
        \mathrm{hi}_t(c):=\sum_km_k\mathbf 1[c\le\gamma_{k,t}],\qquad\mathrm{lo}_t(c):=\sum_km_k\mathbf 1[c<\gamma_{k,t}].
    \]
    Если $a_t\in S'_k$, то по леммам~\ref{lem:interior} и~\ref{lem:ray} $e^{-\rho}\le\gamma_{k,t}$. Если $b_t\notin S'_k$, то $b_t\in K_k$ и $\gamma_{k,t}\le e^{\rho}$. Поэтому
    \[
        \mathrm{hi}_t(e^{-\rho})\ge w(a_t)\ge y_t\ge w(b_t)\ge\mathrm{lo}_t(e^{\rho}).
    \]
    Функция $\mathrm{hi}_t$ не возрастает, и множество $\set{c\mid\mathrm{hi}_t(c)\ge y_t}$ замкнуто. Положим $c_t:=\max\set{c\in[e^{-\rho},e^{\rho}]\mid\mathrm{hi}_t(c)\ge y_t}$. Это множество содержит $e^{-\rho}$, поэтому $c_t$ определено и $\mathrm{hi}_t(c_t)\ge y_t$. Покажем $\mathrm{lo}_t(c_t)\le y_t$. Если $c_t=e^{\rho}$, это доказано выше. Иначе $\mathrm{hi}_t(c)<y_t$ при $c\in(c_t,e^{\rho}]$, и $\mathrm{lo}_t(c_t)=\lim_{c\downarrow c_t}\mathrm{hi}_t(c)\le y_t$.

    Положим $q_t:=c_tp_t$ и $A_k:=\set{t\mid c_t<\gamma_{k,t}}$, $B_k:=\set{t\mid c_t\le\gamma_{k,t}}$. Тогда $\mathrm{lo}=\sum_km_k\mathbf 1_{A_k}\le y\le\sum_km_k\mathbf 1_{B_k}=\mathrm{hi}$ покоординатно. По лемме~\ref{lem:cube} ниже $y=\sum_j\mu_j\mathbf 1_{C_j}$, $\mu_j>0$, где каждое $C_j$ удовлетворяет $A_k\subset C_j\subset B_k$ при некотором $k$. Осталось проверить $C_j\in\widehat{\Sp}(Q)$, и тогда $y$ слабо разрешим при $Q$ по теореме~\ref{thm:closure}.

    Пусть $A_k\subset C\subset B_k$ и $s\notin C$. Тогда $s\notin A_k$, то есть $c_s\ge\gamma_{k,s}$, и $q_s\in K_k$ по лемме~\ref{lem:ray}. Поэтому $K_C(Q)\subset K_k$. Для $t\in C\subset B_k$ имеем $c_t\le\gamma_{k,t}$, и $q_t\notin\interior K_k\supset\interior K_C(Q)$ по лемме~\ref{lem:ray}. По лемме~\ref{lem:interior} $C\in\widehat{\Sp}(Q)$.
\end{proof}

\begin{lemma}[куб]\label{lem:cube}
    Пусть $m_k>0$, $A_k\subset B_k\subset[T]$ при $k\in[n]$, и $\sum_km_k\mathbf 1_{A_k}\le y\le\sum_km_k\mathbf 1_{B_k}$. Тогда $y=\sum_j\mu_j\mathbf 1_{C_j}$ с $\mu_j>0$, причём для каждого $j$ найдётся $k$ с $A_k\subset C_j\subset B_k$.
\end{lemma}

\begin{proof}
    Разместим на $[0,M)$, $M:=\sum_km_k$, последовательные интервалы $I_k$ длины $m_k$ и обозначим через $k(x)$ номер интервала, содержащего $x$. Для $t\in[T]$ положим $E_t:=\bigcup_{k:\,t\in B_k\setminus A_k}I_k$ и $\ell_t:=\sum_km_k\mathbf 1_{A_k}(t)$. Тогда $\abs{E_t}=\sum_km_k\bigl(\mathbf 1_{B_k}(t)-\mathbf 1_{A_k}(t)\bigr)\ge y_t-\ell_t\ge0$. Возьмём начальный отрезок $F_t:=E_t\cap[0,x_t)$ меры $y_t-\ell_t$ и положим $C(x):=A_{k(x)}\cup\set{t\mid x\in F_t}$. Тогда $A_{k(x)}\subset C(x)\subset B_{k(x)}$, мера множества $\set{x\mid t\in C(x)}$ равна $\ell_t+(y_t-\ell_t)=y_t$, а $C(x)$ кусочно постоянна с конечным числом интервалов постоянства.
\end{proof}

Лемму можно получить и из \lean{PolyHardness.cube_mem_span} (\texttt{poly-hardness/}).

\section{Точки излома}

Из определения $\widehat{\Sp}$ видно, что оно не меняется при общем положительном масштабе точек: $\widehat{\Sp}(\lambda W)=\widehat{\Sp}(W)$ при $\lambda>0$. Поэтому семейство для $W=(e^{-\rho}P,e^{\rho}P)$ совпадает с семейством для $(P,zP)$, $z:=e^{2\rho}$, и индексы те же. Условие (C) зависит от $\rho$ только через $z$.

\paragraph{Однородные векторы.} Точке $w_r$ сопоставим вектор $(w_r,1)\in\R^{d+1}$ и добавим векторы $(e_i,0)$, $i\in[d]$. Для линейной функции $\eta=(\xi,\eta_0)$ на $\R^{d+1}$ положим $\eta(w,1)=\ip{\xi}{w}+\eta_0$.

\begin{lemma}[слабая достижимость как условие на знаки]\label{lem:signs}
    Индекс $t\in S$ имеет слабого свидетеля тогда и только тогда, когда существует ненулевая линейная функция $\eta$ на $\R^{d+1}$ с $\eta(w_t,1)=0$, $\eta(w_s,1)\ge0$ при $s\notin S$ и $\eta(e_i,0)\ge0$ при $i\in[d]$.
\end{lemma}

\begin{proof}
    По свидетелю $\xi$ положим $\eta:=(\xi,-\ip{\xi}{w_t})$. Тогда $\eta(w_s,1)=\ip{\xi}{w_s-w_t}$, $\eta(e_i,0)=\xi^i$, и $\eta\ne0$, так как $\xi\ne0$. Обратно, из $\eta(w_t,1)=0$ следует $\eta_0=-\ip{\xi}{w_t}$, и $\eta\ne0$ равносильно $\xi\ne0$.
\end{proof}

Знаковым вектором функции $\eta$ называется набор знаков $\eta(v_j)$ по всем векторам $v_j$ конфигурации.

\begin{lemma}[локальное постоянство]\label{lem:local}
    Пусть $I\subset\R$~--- интервал, $v_j:I\to\R^{d+1}$, $j\in J$,~--- конечный набор непрерывных функций, векторы $v_j(\rho)$ при каждом $\rho$ порождают $\R^{d+1}$. Пусть $\rho_0\in I$ и каждый определитель $\det(v_j)_{j\in B}$, $B\subset J$, $\abs{B}=d+1$, либо тождественно равен нулю на $I$, либо не равен нулю в $\rho_0$. Тогда всякий знаковый вектор ненулевой функции при $\rho_0$ реализуется ненулевой функцией при всех $\rho$ из некоторой окрестности $\rho_0$.
\end{lemma}

\begin{proof}
    Определители, не равные тождественно нулю, не равны нулю в $\rho_0$ и по непрерывности в некоторой окрестности $N\ni\rho_0$. Поэтому в $N$ каждый определитель либо равен нулю тождественно, либо не равен нулю ни в одной точке.

    \emph{Независимость не зависит от $\rho\in N$.} Пусть $Z'\subset J$ независимо при $\rho_1\in N$. Так как $v_j(\rho_1)$ порождают $\R^{d+1}$, достроим $Z'$ до базиса $B\subset J$. Тогда $\det_B(\rho_1)\ne0$, значит, определитель не равен нулю тождественно, и $\det_B(\rho)\ne0$ при всех $\rho\in N$. Поэтому $Z'$ независимо при всех $\rho\in N$. В частности, ранг любого $Z\subset J$ постоянен на $N$.

    Пусть $\eta$~--- ненулевая функция при $\rho_0$, $Z:=\set{j\mid\eta(v_j(\rho_0))=0}$ и $r:=\operatorname{rank}Z$. Выберем независимое $Z_0\subset Z$ с $\abs{Z_0}=r$. Оно независимо и при всех $\rho\in N$, и порождает $\operatorname{span}\set{v_j(\rho)}_{j\in Z}$, так как ранг $Z$ равен $r$. Достроим $Z_0$ набором $C\subset J$, $\abs{C}=d+1-r$, до базиса. Он является базисом при всех $\rho\in N$. Определим $\eta_\rho$ условиями
    \[
        \eta_\rho(v_j(\rho))=0\ (j\in Z_0),\qquad\eta_\rho(v_c(\rho))=\eta(v_c(\rho_0))\ (c\in C).
    \]
    Коэффициенты $\eta_\rho$ непрерывны по $\rho$, и $\eta_{\rho_0}=\eta$. Для $j\in Z$ вектор $v_j(\rho)$ лежит в $\operatorname{span}Z_0(\rho)$, поэтому $\eta_\rho(v_j(\rho))=0$ при всех $\rho\in N$. Для $j\notin Z$ значение $\eta_\rho(v_j(\rho))$ непрерывно и в $\rho_0$ отлично от нуля, поэтому знак сохраняется вблизи $\rho_0$. Индексов конечное число, поэтому окрестность можно выбрать общей. Функция $\eta_\rho$ ненулевая, так как конфигурация порождает $\R^{d+1}$ и $\eta_\rho(v_j(\rho))\ne0$ при некотором $j$ (знаковый вектор $\eta$ не нулевой).
\end{proof}

\begin{corollary}\label{cor:local}
    Пусть в точке $z_0>0$ каждый определитель конфигурации $(p_t,1)$, $(zp_t,1)$, $(e_i,0)$ либо равен нулю тождественно по $z$, либо не равен нулю при $z_0$. Тогда $\widehat{\Sp}(P,zP)\supseteq\widehat{\Sp}(P,z_0P)$ для всех $z$ из некоторой окрестности $z_0$.
\end{corollary}

\begin{proof}
    Пусть $S\in\widehat{\Sp}(P,z_0P)$ и $t\in S$. По лемме~\ref{lem:signs} функция $\eta$ для пары $(S,t)$ задаёт знаковый вектор с нулём на $v_t$ и неотрицательными знаками на $v_s$, $s\notin S$, и на $(e_i,0)$. По лемме~\ref{lem:local} такой вектор реализуется ненулевой функцией при всех $z$ вблизи $z_0$. Пар $(S,t)$ конечное число.
\end{proof}

\paragraph{Рациональность.}
Рассмотрим $(d+1)$-подмножество $B$ конфигурации $(P,zP)$ с $\alpha$ строками вида $(p_t,1)$, $\beta$ строками вида $(zp_s,1)$ и остальными вида $(e_i,0)$. Разделим строки второго вида на $z>0$. Они станут $(p_s,1/z)$, и определитель домножится на $z^{-\beta}$. Первые $d$ столбцов полученной матрицы $X$ не зависят от $z$, а последний столбец равен $\mathbf u+z^{-1}\mathbf v$, где $\mathbf u$ и $\mathbf v$~--- индикаторы строк первого и второго вида. Определитель линеен по последнему столбцу, поэтому
\[
    \det{}_B(z)=z^{\beta}\bigl(D_0+z^{-1}D_1\bigr)=z^{\beta-1}(D_0z+D_1),\qquad D_0=\det[X\mid\mathbf u],\ D_1=\det[X\mid\mathbf v].
\]
Если $D_0=0$, то определитель равен нулю тождественно или не обращается в нуль при $z>0$. Если $D_0\ne0$, то у него единственный корень $z=-D_1/D_0$. Если $P$ рационально, то $D_0,D_1\in\Q$. \textit{Точками излома} назовём корни $z>1$ определителей, не равных нулю тождественно. Они рациональны. Подмножеств $B$ не более $\binom{2T+d}{d+1}$, поэтому точек излома $N=O\bigl((2T+d)^{d+1}\bigr)$.

\section{Теорема о наименьшем радиусе}

Пусть $F:=\set{\rho\ge0\mid\text{выполнено (A)}}$.

\begin{theorem}[наименьший радиус]\label{thm:least}
    Если $F\ne\varnothing$, то $F=[\rho^*,\infty)$. При этом $\rho^*=0$ или $z^*:=e^{2\rho^*}$ является точкой излома.
\end{theorem}

\begin{proof}
    \emph{Монотонность.} Шары вложены друг в друга, поэтому из $\rho\in F$ следует $[\rho,\infty)\subset F$.

    \emph{Замкнутость.} Пусть $\rho_n\in F$, $\rho_n\downarrow\rho$. По теореме~\ref{thm:corners} найдутся $c^{(n)}\in[e^{-\rho_n},e^{\rho_n}]^T$, при которых $y$ слабо разрешим при $c^{(n)}P$. Переходя к подпоследовательности, считаем $c^{(n)}\to c\in[e^{-\rho},e^{\rho}]^T$. Слабая разрешимость замкнута. Действительно, любая окрестность $U$ точки $cP$ содержит $c^{(n)}P$ при большом $n$ вместе с некоторой окрестностью $c^{(n)}P$, и в ней по определению есть набор, при котором $y$ разрешим. Значит, $y$ слабо разрешим при $cP$, и $\rho\in F$ по теореме~\ref{thm:corners}.

    \emph{Точка излома.} Пусть $\rho^*>0$ и $z^*$ не является точкой излома. Тогда каждый определитель конфигурации $(P,zP)$ либо равен нулю тождественно, либо не равен нулю при $z^*$, так как по предыдущему разделу его корень при $z>1$ единственный и $z^*>1$ им не является. По теореме~\ref{thm:corners} выполнено (C) при $\rho^*$ для $W^*=(e^{-\rho^*}P,e^{\rho^*}P)$. По следствию~\ref{cor:local} при $z=e^{2\rho}$, $\rho<\rho^*$ достаточно близком к $\rho^*$, семейство $\widehat{\Sp}(W)$ содержит $\widehat{\Sp}(W^*)$. Индексы $a_t,b_t$ и значения $y_t$ те же, поэтому $w$ из (C) при $\rho^*$ годится и при $\rho$. По теореме~\ref{thm:corners} $\rho\in F$, что противоречит определению $\rho^*$.
\end{proof}

\begin{corollary}\label{cor:gap}
    На каждом интервале $(z_i,z_{i+1})$ между соседними точками излома (и на $(1,z_1)$ и $(z_N,\infty)$) допустимость постоянна.
\end{corollary}

\section{Алгоритм}

\paragraph{Оракул.} Для набора точек $W$ с индексами $R$ и границ $\mathrm{lo},\mathrm{hi}\in[-\infty,\infty]^R$ оракул отвечает, существует ли $w\in\cone\set{\mathbf 1_S\mid S\in\widehat{\Sp}(W)}$ с $\mathrm{lo}\le w\le\mathrm{hi}$. Условие (C) есть запрос с $W=(e^{-\rho}P,e^{\rho}P)$, $\mathrm{lo}=(y,0)$, $\mathrm{hi}=(\infty,y)$ по индексам $(a,b)$. При $\mathrm{lo}=\mathrm{hi}=y$ получаем слабую разрешимость по теореме~\ref{thm:closure}.

\paragraph{Шаги.} Пусть $1<z_1<\dots<z_N$~--- точки излома, $z_0:=1$, $z_{N+1}:=\infty$.
\begin{enumerate}
    \item Проверить $z=1$ (запрос с $\mathrm{lo}=\mathrm{hi}=y$ при $W=P$). Если допустимо, то $\rho^*=0$.
    \item Для $i\in[N]$ выбрать рациональное $k_i$ с $k_i^2\in(z_i,z_{i+1})$. Запрос допустим при $W=(P/k_i,Pk_i)$ тогда и только тогда, когда допустим весь интервал (следствие~\ref{cor:gap}). Если интервал $(z_N,\infty)$ недопустим, то $F=\varnothing$.
    \item По теореме~\ref{thm:least} $F=[z^*,\infty)$ в переменной $z$ и $z^*\in\set{1,z_1,\dots,z_N}$. Если $z=1$ недопустимо, то $z^*>1$, и $z^*=z_i$ для наименьшего $i\ge1$, при котором интервал $(z_i,z_{i+1})$ допустим. Его находим бинарным поиском. Тогда $\rho^*=\tfrac12\ln z_i$.
\end{enumerate}
Число обращений к оракулу равно $O(\log N)=O(d\log(T+d))$. Оно не зависит от длины записи цен. Точки излома рациональны, поэтому $k_i$ и все данные запросов рациональны. Перечисление точек излома требует $\binom{2T+d}{d+1}$ определителей порядка $d+1$ на числах длины $O(dL)$, где $L$~--- длина записи цен.

\begin{remark}[бисекция без перечисления]
    При больших $d$ перечисление дорого. Например, при $d=7$, $T=40$ определителей около $6\cdot10^{10}$. Тогда достаточно бисекции по $z$. Умножим каждую строку на положительное число так, чтобы она стала целой. Отношение $z=-D_1/D_0$ при этом не меняется. Пусть $B$~--- наибольший модуль полученных координат. По неравенству Адамара $\abs{D_0},\abs{D_1}\le H:=(d+1)^{(d+1)/2}B^d$. Поэтому $z^*\le H$, а две различные точки излома отстоят не меньше чем на $1/H^2$. После $3\log_2H=O(dL)$ шагов бисекции на отрезке длины меньше $1/H^2$ остаётся не больше одной точки излома. Её находят цепной дробью без обращений к оракулу. Число обращений зависит от $L$, то есть алгоритм слабо полиномиален по числу обращений.
\end{remark}

\todo{Для статьи: основной алгоритм~--- поиск по точкам излома (теорема~\ref{thm:least}, число обращений не зависит от $L$). Бисекцию дать одним замечанием. Разворачивать её, только если в статье будут расчёты при $d\ge4$. Сравнение сильной и слабой полиномиальности и теорему Тардош подробно не излагать.}

\paragraph{Сертификаты.} При ответе YES оракул выдаёт столбцы $S'_k$ и массы $m_k$. По ним строятся $c_t$ из доказательства (C)$\Rightarrow$(B) и $Q=(c_tp_t)$, и точной арифметикой проверяются $\dd(q_t,p_t)\le\rho$ и $y=\sum_j\mu_j\mathbf 1_{C_j}$ с $C_j\in\widehat{\Sp}(Q)$. При ответе NO оракул выдаёт вектор $\lambda\in\R^R$ с $\lambda(S):=\sum_{r\in S}\lambda_r\ge0$ для всех $S\in\widehat{\Sp}(W)$ и
\[
    \sum_r\lambda_r^-\,\mathrm{lo}_r>\sum_r\lambda_r^+\,\mathrm{hi}_r.
\]
Это лемма Фаркаша для системы $m\ge0$, $\mathrm{lo}\le\sum_Sm_S\mathbf 1_S\le\mathrm{hi}$. Проверка условия $\lambda(S)\ge0$ на всём $\widehat{\Sp}(W)$ требует оракула столбцов. В сертификате $\rho^*$ оба ответа берутся на соседних интервалах $(z_{i-1},z_i)$ и $(z_i,z_{i+1})$.

\paragraph{Сложность оракула.} При $d=2$ запрос есть одна задача линейного программирования о потоке в графе цепочек (\texttt{poly-hardness/}), и весь алгоритм полиномиален по $T$. При $d\ge3$ задача NP-трудна (\texttt{np-hardness/}).

\paragraph{Генерация столбцов.} На практике оракул реализуется генерацией столбцов. Главная задача есть задача линейного программирования по уже найденным слабо достижимым множествам. Её двойственные переменные задают веса $g\in\R^R$. Вспомогательная задача ищет слабо достижимое $S$ с наибольшим $g(S)$. Это MILP с одним множеством: бинарные $x_r$ при $r\in R$ и свидетели $\xi_t$ в симплексе. Ограничения свидетелей
\[
    \sum_{i\in[d]}\xi_t^i\,\frac{w_s^i-w_t^i}{M_{st}}\ge-(1-x_t+x_s),\qquad M_{st}:=\max_i\abs{w_s^i-w_t^i}.
\]
Так как $\xi_t$ лежит в симплексе, при $x_t=0$ или $x_s=1$ ограничение выполнено всегда, а при $x_t=1$, $x_s=0$ оно равносильно $\ip{\xi_t}{w_s-w_t}\ge0$. Поэтому постановка точна. Деление на $M_{st}$ существенно. Без него решатель с допусками порядка $10^{-6}$ принимает множества, которые не слабо достижимы. Каждый найденный столбец проверяется точно.

При $d=2$ задача о потоке есть компактная запись той же задачи, а генерация столбцов есть её запись по путям (разложение Данцига--Вулфа). Вспомогательная задача при $d=2$ есть задача о длиннейшем пути в графе цепочек. Она решается точно за $O(T^4)$.

\paragraph{Точность.} Постановка MILP точна, но решатель работает в плавающей арифметике. Ответ YES всегда точен, так как столбцы и массы проверяются в рациональных числах. Ответ NO опирается на оптимальность вспомогательной задачи. При $d=2$ он подтверждается одним длиннейшим путём для вектора $\lambda$ на интервале под $z^*$. При $d\ge3$ для точности нужен решатель с рациональной арифметикой, его достаточно вызвать один раз на сертификате. Допуск оптимальности решателя должен быть нулевым.

\todo{Для статьи: граф цепочек нужен для теоремы о полиномиальности при $d=2$ и для итогового сертификата NO. Вычисления ведутся генерацией столбцов. Цикл повторных решений с отсечением отвергнутых множеств при $d=2$ не срабатывает, а при $d=7$ срабатывает и не спасает. Описывать стоит точную вспомогательную задачу, а не этот цикл.}

\section{Пример}

Данные EU KLEMS, США, вся экономика, 1995--2020, $T=26$, $d=2$. Цены есть цены капитала и труда, нормированные ценой промежуточных затрат, выпуск $y=\mathrm{GO\_Q}$. Результат: $\rho^*=0.089654$, точек излома 8310, обращений к оракулу 15, оба сертификата проверены точно.

Для $d=3$ к ценам добавляется цена выпуска $\mathrm{GO\_CP}/\mathrm{GO\_Q}$, где $\mathrm{GO\_CP}=\mathrm{VA\_CP}+\mathrm{II\_CP}$. Точек излома 127160, они вычисляются за 30 секунд.

\paragraph{Отрасли США.} Данные BEA--BLS (интегрированный отраслевой счёт производства, исследовательское продолжение на 1947--2016 годы). Это 63 отрасли за 1963--2016 годы и 37 отраслей за 1947--2016 годы после склейки по 1963 году. Цены капитала и труда агрегированы индексом Торнквиста. Точки $(p_K/p_{II},p_L/p_{II})$, выпуск $y=\mathrm{GO\_Q}$, $d=2$.

При $T=40$ (1977--2016) во всех 63 отраслях $y$ не слабо разрешим при исходных ценах, $\rho^*$ от $0.02$ до $0.89$. Точек излома около $3\cdot10^4$, обращений к оракулу 16--17. Оракул на потоке и генерация столбцов дают одинаковые $z^*$ во всех отраслях. Суммарное время 2350 с у потока и 960 с у генерации столбцов. Генерация столбцов быстрее в 45 отраслях из 63.

С ростом $T$ разрыв растёт в пользу генерации столбцов.
\begin{center}
\begin{tabular}{lrrrr}
    \toprule
    отрасль & $T=40$ & $T=50$ & $T=60$ & $T=70$ \\
    \midrule
    Machinery & 126 / 22 & 471 / 64 & 1040 / 72 & 2339 / 197 \\
    Rental and leasing & 100 / 43 & 161 / 57 & 447 / 135 & 1134 / 286 \\
    Utilities & 4.6 / 13.9 & 17 / 22 & 72 / 35 & 170 / 94 \\
    Petroleum and coal & 6.2 / 16.6 & 15 / 24 & 28 / 27 & 59 / 50 \\
    \bottomrule
\end{tabular}
\end{center}
В таблице время поиска $\rho^*$ в секундах: поток / генерация столбцов. При $T=60$ точек излома около $10^5$, при $T=70$ около $1.6\cdot10^5$, обращений 19--20. Время построения графа цепочек в точной арифметике входит во время потока и растёт как $O(T^4)$. При $T=40$ это около 5 с на обращение.

\paragraph{Синтетика.} Чтобы найти примеры, где оракул на потоке быстрее, рассмотрены пять синтетических семейств при $d=2$. Это случайные точки (A), точки около вогнутой дуги (B) и выпуклой дуги (C), точки целочисленной решётки (D) и вогнутая дуга со случайными $y$ (E). На вогнутой дуге каждая точка лежит выше хорд соседей, и слабо достижимы только множества вида отрезков. На выпуклой дуге слабо достижимо любое множество. Решётка даёт много совпадающих координат. Измерялось одно обращение к оракулу при $k=11/10$.
\begin{center}
\begin{tabular}{lrrrrr}
    \toprule
    семейство & $T=20$ & $T=30$ & $T=40$ & $T=50$ & $T=60$ \\
    \midrule
    A & \textbf{0.17} / 0.19 & \textbf{0.96} / 1.55 & 1.24 / 0.51 & 2.7 / 2.0 & 13.4 / 4.5 \\
    B & \textbf{0.19} / 0.51 & 0.96 / 0.62 & 2.6 / 2.3 & 6.8 / 1.2 & 15.3 / 6.0 \\
    C & 0.68 / 0.34 & 3.5 / 0.85 & 14.9 / 1.6 & 57.7 / 2.8 & 134.7 / 3.9 \\
    D & \textbf{0.07} / 0.13 & \textbf{0.43} / 0.84 & 1.18 / 1.22 & 2.7 / 2.5 & 6.7 / 4.0 \\
    E & \textbf{0.22} / 0.33 & 0.92 / 0.71 & 2.7 / 1.6 & 5.9 / 3.4 & 13.9 / 8.5 \\
    \bottomrule
\end{tabular}
\end{center}
В таблице время в секундах: поток / генерация столбцов. Жирным выделены случаи, где поток быстрее. Ответы двух оракулов совпали во всех 25 случаях. Поток быстрее только при $T\le30$: на случайных точках, на вогнутой дуге и на решётке, до 2.7 раза. Начиная с $T=40$ генерация столбцов быстрее или наравне во всех семействах. Худший случай для потока~--- выпуклая дуга, где граф цепочек наибольший. При $T=60$ разрыв 35 раз.

Это только синтетика. Семейства подобраны вручную, и они не исчерпывают возможных трудных примеров. Пример, где поток быстрее при больших $T$, найти не удалось. Это не доказывает, что такого примера нет.

\paragraph{Большие $d$.} При $d\ge3$ оракул NP-труден, а точек излома $\binom{2T+d}{d+1}$. Поэтому $\rho^*$ находится бисекцией по $\rho$ с точностью $10^{-5}$ (замечание о бисекции), оракул~--- генерация столбцов. Наборы данных следующие.
\begin{itemize}
    \item EU KLEMS, США, 1995--2020, $d=3$: цены капитала, труда и промежуточных затрат, нормированные ценой выпуска $\mathrm{GO\_CP}/\mathrm{GO\_Q}$, $\mathrm{GO\_CP}=\mathrm{VA\_CP}+\mathrm{II\_CP}$.
    \item Berndt--Wood, обрабатывающая промышленность США, 1947--1971, $d=4$: цены капитала, труда, энергии и материалов, нормированные ценой выпуска.
    \item BEA--BLS, 1977--2016, $d=4$: агрегированный капитал, труд с высшим образованием и без, промежуточные затраты, нормированные ценой выпуска.
    \item BEA--BLS, 1977--2016, $d=7$: пять видов капитала и два вида труда, нормированные ценой промежуточных затрат.
\end{itemize}
\begin{center}
\begin{tabular}{llrrrrr}
    \toprule
    данные & $d$, $T$ & $\rho^*$ & обращений & время, с & макс.\ обращение, с & итераций (макс.) \\
    \midrule
    EU KLEMS, США & 3, 26 & 0.03461 & 19 & 30 & 3.1 & 8 \\
    Berndt--Wood & 4, 25 & 0.01600 & 19 & 18 & 2.3 & 16 \\
    BEA, Machinery & 4, 40 & 0.03843 & 19 & 80 & 13.5 & 18 \\
    BEA, Petroleum and coal & 4, 40 & 0.13562 & 19 & 23 & 2.6 & 7 \\
    BEA, Rental and leasing & 4, 40 & 0.01656 & 19 & 80 & 26.5 & 34 \\
    BEA, Utilities & 4, 40 & 0.04676 & 19 & 128 & 38.9 & 39 \\
    BEA, Machinery & 7, 40 & 0.03165 & 19 & 299 & 68.7 & 23 \\
    BEA, Utilities & 7, 40 & 0.17403 & 19 & 267 & 39.8 & 24 \\
    BEA, Petroleum and coal & 7, 40 & \multicolumn{4}{l}{сбой на 18-м обращении} \\
    BEA, Rental and leasing & 7, 40 & \multicolumn{4}{l}{сбой на 10-м обращении} \\
    \bottomrule
\end{tabular}
\end{center}
Значения $\rho^*$ даны с точностью $10^{-5}$. Прежний MILP с $K=2T$ множествами на данных EU KLEMS при $d=3$ не дал ответа ни на одно обращение за 300 с. Генерация столбцов решает случаи с $d\le4$ за минуты. При $d=7$ в двух отраслях из четырёх вспомогательная задача 20 раз подряд вернула множества, которые не прошли точную проверку. Это численная проблема ограничений со свидетелями при большом $d$. Подзадача остаётся NP-трудной, и эти примеры не дают гарантии для других данных.

\todo{При $d=7$ нужна вспомогательная задача без непрерывных свидетелей: ленивые отсечения $x_t\le\sum_{s\in A}x_s$, $\abs{A}\le d+1$, с точной геометрией. Цикл повторных решений с отсечением отвергнутых множеств при $d=7$ срабатывает, то есть он не избыточен.}

\section{Lean}

Библиотека \lean{HexPerturbation}, сборка \texttt{lake build HexPerturbation}.

\begin{center}
\begin{tabular}{p{0.38\textwidth}p{0.52\textwidth}}
    \toprule
    Утверждение & Lean \\
    \midrule
    лемма~\ref{lem:metric} & \lean{abs_log_sub_le_projDist}, \lean{projDist_smul_le} \\
    теорема~\ref{thm:corners} (углы) & \lean{proj_iff_scale}, \lean{proj_iff_corners}; при $d=2$ \lean{hex_iff_scale}, \lean{hex_iff_corners} \\
    лемма~\ref{lem:signs} & \lean{witness_iff_homogeneous} \\
    лемма~\ref{lem:local} & \lean{eventually_exists_sign_eq} \\
    следствие~\ref{cor:local} & \lean{eventually_isWeakReachable} \\
    теорема~\ref{thm:least} (наименьший радиус) & \lean{least_radius_breakpoint} \\
    лемма~\ref{lem:shape} (форма шара) & \lean{projDist_le_iff_mem_convexHull} \\
    формула $\dd$ при $d=2$ & \lean{hexDist_eq_projDist} \\
    рациональность $e^{2\rho^*}$ & \lean{least_radius_rational} \\
    \bottomrule
\end{tabular}
\end{center}

Не формализованы алгоритм и его сложность (в Mathlib нет классов сложности) и проверки сертификатов в коде. Их корректность опирается на формализованные теоремы.

\end{document}
